\section[5. GUI Programming - {\it 그래픽 프로그래밍}]{5. GUI Programming}

\subsection{TUI vs GUI}

\textbf{텍스트 유저 인터페이스}\textit{\textsuperscript{Text User Interface, TUI}}, \textbf{그래픽 유저 인터페이스}\textit{\textsuperscript{Graphic User Interface, GUI}}는 사용자-컴퓨터 간의 상호작용 위한 주요한 인터페이스 유형이다.\\
TUI는 텍스트 기반이다. 콘솔이나 명령 줄로 알려져 있음. 사용자는 명령어를 입력하고 텍스트 기반 정보를 받아들이며 키보드 이용하여 상호작용한다. MS DOS, linux shell, vim 등이 그 예시이다.
\begin{itemize}
\item 장점: 자원 소모가 적고 경량. 텍스트 기반이라 console에서도 사용 가능. 간단하고 빠르며, 스크립트와의 통합이 용이함.
\item 단점: 사용자가 명령어나 옵션 기억하고 입력해야 함. 그래픽 요소 제한적이기에 복잡한 시각적 표현 어려움
\end{itemize}

GUI는 그래픽 요소를 사용하여 사용자와 상호작용한다. 사용자는 윈도우, 버튼, 메뉴, 입력 필드와 같은 요소를 마우스로 조작한다.
\begin{itemize}
\item 장점: 시각적인 요소로 직관적인 환경을 제공하며, 복잡한 작업을 그래픽으로 간단하게 수행함. 다양한 입력 방식 지원하여 사용자 편의성 높임.
\item 단점: 자원 소모 크고, 구현과 디자인에 시간이 많이 걸림. 다양한 플랫폼에서 동일한 모습을 유지하기가 힘듦.
\end{itemize}

\subsection{Tkinter}

파이썬은 GUI 프로그램을 제작하는데 쓸 수 있는 tkinter, pyqt, pyside, kivy, wxpython 등의 다양한 도구 지원한다. 이중 tkinter는 파이썬의 표준 라이브러리로 문법이 상대적으로 단순하며 별도의 설치가 필요 없다는 장점을 갖는다.\\
또한, tkinter는 다음의 특징을 갖는다.
\begin{itemize}
\item \textbf{크로스 플랫폼 지원}: 파이썬은 크로스 플랫폼 개발을 지원하므로 동일한 코드를 사용하여 여러 플랫폼에서 GUI 애플리케이션을 실행할 수 있다.
\item \textbf{확장성}: 파이썬은 다른 언어와의 통합이 용이하며, C/C++로 작성된 라이브러리의 결합을 통해 추가적인 기능을 제공할 수 있다. PyQT, wxPython은 Qt, wxWidgets라는 C++ 기반 프레임워크 이용할 수 있다.
\item \textbf{풍부한 기능과 위젯}: 다양한 기능과 위젯을 제공하여 사용자 인터페이스를 다양하게 구성할 수 있음. 버튼, 레이블, 텍스트 상자, 체크박스, 라디오 버튼, 드롭다운 메뉴 등의 위젯 활용할 수 있음. 그래픽 요소도 처리할 수 있기 때문에, 도표, 차트, 이미지 표시와 같은 고급 기능을 개발하는데 유용함.
\item \textbf{데이터 시각화}: matplotlib, seaborn, plotly 등을 이용하여 데이터 시각적으로 표현하고 그래프, 플롯, 차트 등을 생성할 수 있음. 데이터 시각화는 데이터 분석 및 보고에 필수적인 요소로 사용이 가능함.
\end{itemize}

\textbf{위젯}\textit{\textsuperscript{Widget}}은 사용자 인터페이스를 구성하는 요소를 통칭하는 말이다. 위젯은 다양한 상호작용 가능한 요소를 포함한다. 즉, 위젯은 사용자와 애플리케이션 간의 상호작용을 담당하며 동시에 애플리케이션 기능 제어하고 정보를 표시하는 역할을 한다.
\begin{itemize}
\item label: 사용자에게 텍스트나 이미지 표현
\item button: 사용자가 클릭할 수 있는 버튼
\item entry: 사용자로부터 텍스트를 입력받을 수 있는 상자
\item checkbutton: 사용자가 체크할 수 있는 옵션을 표시. 여러 개의 체크 박스를 그룹으로 묶어 다중 선택 가능함
\item radiobutton: 여러 개의 옵션 중 하나를 선택할 수 있는 버튼. 여러 개의 라디오 버튼을 그룹으로 묶어 단일 선택 가능함
\item listbox: 여러 항목을 리스트 형태로 표시하고 사용자가 선택할 수 있게 하며, 단일 선택과 다중 선택이 가능함
\item combobox: 텍스트 상자와 드롭다운 목록을 결합한 widget. 사용자가 선택할 수 있는 목록을 제공하며, 자유롭게 텍스트 입력 가능함
\item text: 여러 줄의 텍스트를 입력하거나 표시하는데 사용. 대용량 텍스트 편집 가능하며 스크롤바 포함 가능함.
\end{itemize}

다음은 Tkinter를 사용해 GUI 프로그램을 작성할 때 공통적으로 사용하는 명령어들이다. 즉, 최소한 아래의 명령어를 작성해야 Tkinter를 사용해 프로그래밍을 시작할 준비가 된 것이다.
\begin{itemize}
\item window=tk.Tk(): Tk class로 주 창을 생성. 루트 윈도우를 초기화.
\item window.title(): 창 제목 설정
\item window.mainloop(): tkinter 애플리케이션의 이벤트 루프를 실행한다. 이벤트 처리하고 창을 업데이트하는데 사용된다.
\end{itemize}

위젯을 만들고 프로그램의 흐름을 제어하는 데에는 아래의 명령어를 주로 사용한다.
\begin{itemize}
\item 위젯 = tk.Entry(): 입력 창 설정. input\_entry.get()으로 입력된 값 받아옴.
\item 위젯.pack(): 화면에 배치. 자동으로 배치됨.
\item 위젯.config(): 속성을 변경.
\item canvas.bind("<키이름>", lambda e: 이벤트): canvas에서 발생하는 키 입력을 bind한다.
\end{itemize}

각 위젯마다 속성이 다르고 이를 어떻게 활용하느냐에 따라 다른 형태를 갖기 때문에 이를 잘 암기하고 활용하는 것은 매우 중요하다.\\
추가적으로 outline, create\_oval, create\_rectangle, create\_triangle는 각각 선, 타원, 직사각형, 삼각형을 화면 상에 그리는 함수로 테두리 색상(outline), 내부 색상(fill), 테두리 두께(width) 등의 속성을 갖는다.
